\section{ERR-008: Git Machine Import Errors Not Translated}
\label{sec:err-008}
\subsection*{Metadata}
\begin{tabular}{ll}
\textbf{Issue ID} & ERR-008 \\
\textbf{Severity} & MEDIUM \\
\textbf{Category} & Error Handling \\
\textbf{Branch} & \branchref{fix/ERR-008-git-machine-import-errors}
\textbf{Changes} & \branchdiff{fix/ERR-008-git-machine-import-errors} \\
\end{tabular}

\subsection*{Executive Summary}
In \srcfile{src/sandbox.ts}, when the optional \code{gitmachine} package fails to load, the catch block throws a hardcoded error message telling the user to install it. However, the actual import error is discarded entirely. If the user has \code{gitmachine} installed but it fails to load due to a broken dependency, version mismatch, or syntax error, they get the misleading message ``Sandbox mode requires the 'gitmachine' package'' instead of the actual error.

\subsection*{Root Cause Analysis}
The catch block discards the real error:

\begin{lstlisting}[language=TypeScript]
// BEFORE — original error discarded
let gitmachine: any;
try {
    gitmachine = await import("gitmachine");
} catch {
    throw new Error(
        "Sandbox mode requires the 'gitmachine' package.\n" +
        "Install it with: npm install gitmachine",
    );
}
\end{lstlisting}

Real scenarios where this causes confusion: the package is installed but has a broken dependency (user sees ``Install gitmachine''), version incompatibility, or a syntax error in the package itself.

\subsection*{Fix Implementation}
The original error is now preserved in the thrown message:

\begin{lstlisting}[language=TypeScript]
// AFTER — original error included
try {
    gitmachine = await import("gitmachine");
} catch (err) {
    throw new Error(
        "Sandbox mode requires the 'gitmachine' package.\n" +
        `Install it with: npm install gitmachine\nOriginal error: ${err instanceof Error ? err.message : String(err)}`,
    );
}
\end{lstlisting}

\subsection*{Affected Files}
\begin{itemize}
    \item \srcfile{src/sandbox.ts} — Lines 51-58
    \item \srcfile{src/\_\_tests\_\_/git-machine-errors.test.ts}
\end{itemize}

\subsection*{Verification}
Tests verify the original error is preserved in the thrown message:

\begin{lstlisting}[language=TypeScript]
const originalErr = new Error("MODULE_NOT_FOUND: cannot find module 'gitmachine'");
try {
    throw new Error(
        "Sandbox mode requires the 'gitmachine' package.\n" +
        `Install it with: npm install gitmachine\nOriginal error: ${originalErr.message}`,
    );
} catch (err: any) {
    assert.ok(err.message.includes("Original error"));
    assert.ok(err.message.includes("MODULE_NOT_FOUND"));
}
\end{lstlisting}
